\Qhead{3.5}{Нелинейные преобразования узкополосного шума}{Нелинейные преобразования (умножения частоты и амплитудное детектирование узкополосного шума).}
\Subfile{3.5.1}{sections/03_statistics/q05/s01}

\recap{Нелинейность порождает суммы и разности частот, а фильтр выбирает измеряемый результат. Гауссовость входа не сохраняется; для степеней комплексной огибающей высшие корреляции можно вычислить через спаривания.}
\board{$x^2$ с медленной и удвоенной составляющими; $n!K^n$; различие $R$ и $R^2$.}
\oral{Почему амплитудный и квадратичный детекторы одного шума дают разные средние и дисперсии?}
\bait{Можно оставить зависимость уширения от формы корреляции: одинаковый умножитель частоты по-разному меняет гауссов и лоренцев спектры.}
\sources{Биленко и соавт.~\cite{bilenko}, §§3.4.4--3.4.7: принципы детектирования. Статистические формулы получены здесь из явно принятой круговой гауссовской модели; это самостоятельное дополнение, а не приписанный пособию расчёт.}
